%% --- thk-thesis: packages that must be loaded after hyperref ---------------
\usepackage[toc,symbols,acronyms,nopostdot,nogroupskip]{glossaries}
\setacronymstyle{long-short}
% No external indexing tool needed: the lists are sorted by LaTeX itself.
\makenoidxglossaries

%% --- Student files (optional, loaded if they exist) -------------------------
\IfFileExists{tex/macros.tex}{%% ---------------------------------------------------------------------------
%% Your own LaTeX definitions. This file is loaded in the preamble.
%% ---------------------------------------------------------------------------

%=== Mathematical operators and symbols ======================================
\DeclareMathOperator{\sgn}{sgn}
\newcommand{\vecW}{\ensuremath{\mathbf{w}}}
\newcommand{\ci}{\ensuremath{\mathrm{i}}}

%=== Hyphenation exceptions ==================================================
% \hyphenation{Ko-ma-Script Bet-ter-BibTeX}

%=== Additional packages ======================================================
% \usepackage{siunitx}
% \usepackage{algorithm}
% \usepackage{algpseudocodex}
}{}
\IfFileExists{tex/glossary.tex}{%% ---------------------------------------------------------------------------
%% Glossary, abbreviations and symbols. Use them in the text with \gls{label}.
%% Only entries that are used in the text appear in the lists.
%% ---------------------------------------------------------------------------

%=== Glossary =================================================================
\newglossaryentry{dog}{
  name={Dog},
  description={Hairy, four-legged mammal. Man's best friend}
}

%=== Abbreviations ============================================================
\newacronym{svm}{SVM}{support vector machine}

%=== Symbols ==================================================================
\newglossaryentry{sym:force}{
  name={\ensuremath{\vec{F}}},
  description={Force, vector quantity},
  sort={F},
  type=symbols,
}
}{}

%% --- thk-thesis: theorem-like statements ------------------------------------
%% Customization 3: Quarto defines theorem, lemma, ... (only those used) above,
%% each with its own counter. Here they share one counter per chapter:
%% Definition 2.1, Theorem 2.2, Lemma 2.3. Algorithms keep their own counter.
\usepackage{amsthm}
\makeatletter
\newcounter{thkstatement}[chapter]
\renewcommand{\thethkstatement}{\thechapter.\arabic{thkstatement}}
\renewcommand{\theHthkstatement}{\thechapter.\arabic{thkstatement}}
\@for\thk@env:={theorem,lemma,corollary,proposition,conjecture,definition,%
    example,exercise,refremark,refsolution}\do{%
  \@ifundefined{c@\thk@env}{}{%
    \expandafter\let\csname c@\thk@env\endcsname\c@thkstatement
    \expandafter\let\csname the\thk@env\endcsname\thethkstatement}}
\makeatother

%% --- thk-thesis: List of Algorithms ------------------------------------------
%% Quarto's algorithms (`::: {#alg-...}`) are statements without a list of
%% their own; thk-thesis.lua writes each one into the file .loa.
\addtotoclist[thk-thesis]{loa}
\setuptoc{loa}{totoc}
\newcommand*{\listofthkalgorithms}{\listoftoc[List of Algorithms]{loa}}
\DeclareTOCStyleEntry[level=1,indent=0pt,numwidth=2.3em,pagenumberwidth=2.4em]{tocline}{algorithm}
